% optimization/alignment/10_continuous_curvature_optimization-DE.tex
\section{Formale Optimierung im Krümmungsraum}

\subsection{Problemstellung}
Es sei \(L>0\) fest und
\[
\kappa:[0,L]\to\R
\]
die Krümmungsfunktion eines Übergangs. Wir nehmen die Randbedingungen
\[
\kappa(0)=0,\qquad\kappa(L)=\kappa_1
\]
und optional
\[
\kappa'(0)=0,\qquad\kappa'(L)=0
\]
an.

\subsection{Zulässige Menge}
\begin{definition}[Zulässige Übergangsmenge]
\[
\mathcal A=
\left\{\kappa\in H^1(0,L)\mid
\kappa(0)=0,\;\kappa(L)=\kappa_1\right\}.
\]
\end{definition}

\subsection{Variationsproblem}
\begin{definition}[Quadratisches Funktional des Krümmungsgradienten]
\[
J[\kappa]=\int_0^L\bigl(\kappa'(s)\bigr)^2\,\dd s.
\]
\end{definition}
\begin{definition}[Optimierungsproblem]
Gesucht ist \(\kappa^\ast\in\mathcal A\) mit
\[
J[\kappa^\ast]=\min_{\kappa\in\mathcal A}J[\kappa].
\]
\end{definition}

\subsection{Euler--Lagrange-Gleichung}
\begin{proposition}
Ist \(\kappa^\ast\) ein Minimierer, dann gilt
\[
\kappa^{\ast\prime\prime}(s)=0.
\]
\end{proposition}

\subsection{Interpretation}
\begin{proposition}
Der Minimierer lautet
\[
\kappa^\ast(s)=\frac{\kappa_1}{L}s.
\]
\end{proposition}

Die Klothoide erscheint somit als Lösung eines Variationsproblems. Dies
verbindet klassischen Eisenbahnübergangsentwurf mit Optimierung im
Krümmungsraum und liefert ein formales Gegenstück zur historischen Rolle der
Klothoide \cite{Glover1900TransitionCurves,Higgins1922TransitionSpiral}.

\subsection{Modell höherer Ordnung}
\begin{definition}
\[
J_2[\kappa]=\int_0^L\bigl(\kappa''(s)\bigr)^2\,\dd s.
\]
\end{definition}
\begin{proposition}
Minimierer erfüllen
\[
\kappa^{(4)}(s)=0.
\]
\end{proposition}
Dies erklärt das Auftreten kubischer Übergangsgesetze und verbindet sie mit
modernem polynomialem Übergangsentwurf
\cite{ZboinskiWoznica2017Polynomial}.

\subsection{Verallgemeinerte Sicht}
Allgemeine Funktionale der Form
\[
J[\kappa]=\int_0^L
F\bigl(\kappa,\kappa',\kappa'';v,c\bigr)\,\dd s
\]
ermöglichen die Einbeziehung dynamischer und ingenieurtechnischer Kriterien.

\subsection{Verbindung zur Hauptthese}
\begin{summary}
Übergangsentwurf kann als Optimierungsproblem im Krümmungsraum formuliert
werden. Klassische Übergangsbögen entstehen als Lösungen solcher Probleme und
stützen die Interpretation von Übergängen als Krümmungssteuerungsfunktionen.
\end{summary}

\section{Gekoppelte Optimierung von Krümmung und Überhöhung}

\subsection{Motivation}
In Eisenbahnanwendungen können Krümmung und Überhöhung nicht unabhängig
optimiert werden. Beide bestimmen gemeinsam die effektive Querbeschleunigung
und deren Änderungsrate.

Dies motiviert ein gekoppeltes Optimierungsproblem mit
\[
\kappa(s)\quad\text{und}\quad\phi(s)
\]
als Entwurfsvariablen. Die Kopplung entspricht Eisenbahnentwurfsnormen, in
denen Krümmung, Geschwindigkeit und Überhöhung keine unabhängigen
Entwurfsgrößen sind \cite{EN13803_2017,Klein1998Trassierung}.

\subsection{Funktional der effektiven Beschleunigung}
\begin{definition}[Effektive Querbeschleunigung]
\[
a_{\mathrm{eff}}(s)=v^2\kappa(s)-g\sin\phi(s).
\]
\end{definition}
Ein natürliches Optimierungsziel ist die Verringerung unausgeglichener
Querbeschleunigung.
\begin{definition}[Beschleunigungsbasiertes Funktional]
Definiere
\[
J_a[\kappa,\phi]
=\int_0^L a_{\mathrm{eff}}(s)^2\,\dd s
=\int_0^L\bigl(v^2\kappa(s)-g\sin\phi(s)\bigr)^2\,\dd s.
\]
\end{definition}

\subsection{Ruckbasiertes Funktional}
\begin{definition}[Ruck]
Der Ruck ist
\[
j(s)=v^3\kappa'(s)-gv\cos\phi(s)\,\phi'(s).
\]
\end{definition}
\begin{definition}[Ruckfunktional]
Definiere
\[
J_j[\kappa,\phi]=\int_0^L j(s)^2\,\dd s,
\]
also
\[
J_j[\kappa,\phi]=\int_0^L
\bigl(v^3\kappa'(s)-gv\cos\phi(s)\,\phi'(s)\bigr)^2\,\dd s.
\]
\end{definition}

Dieses Funktional bestraft schnelle Änderungen der effektiven Beschleunigung
und liefert ein natürliches dynamisches Glattheitskriterium. Solche Kriterien
stehen eng mit Komfort und dynamischer Verträglichkeit in Beziehung
\cite{EN13803_2017,BrustadDalmo2020Review}.

\subsection{Kombinierte Zielfunktion}
\begin{definition}[Gekoppeltes Zielfunktional]
Eine allgemeine gekoppelte Zielfunktion lautet
\[
J[\kappa,\phi]
=\alpha J_a[\kappa,\phi]+\beta J_j[\kappa,\phi]
+\gamma J_\kappa[\kappa]+\delta J_\phi[\phi],
\]
wobei
\[
J_\kappa[\kappa]=\int_0^L(\kappa'(s))^2\,\dd s,
\qquad
J_\phi[\phi]=\int_0^L(\phi'(s))^2\,\dd s.
\]
\end{definition}

Die Gewichte
\[
\alpha,\beta,\gamma,\delta\geq0
\]
bestimmen die relative Bedeutung von Gleichgewicht, Komfort,
Krümmungsglattheit und Überhöhungsglattheit.

\subsection{Zulässige Menge}
\begin{definition}[Zulässige Eisenbahnzustandsmenge]
Es sei
\[
\mathcal B=
\left\{(\kappa,\phi)\;\middle|\;
\kappa\in H^1(0,L),\;\phi\in H^1(0,L),\;
\kappa(0)=0,\;\kappa(L)=\kappa_1\right\}.
\]
\end{definition}

Zusätzliche Randbedingungen können lauten
\[
\phi(0)=0,\qquad\phi(L)=\phi_1,
\]
oder
\[
\kappa'(0)=\kappa'(L)=0,
\qquad\phi'(0)=\phi'(L)=0.
\]

\subsection{Gekoppeltes Optimierungsproblem}
\begin{definition}[Gekoppeltes Eisenbahnoptimierungsproblem]
Gesucht ist
\[
(\kappa^\ast,\phi^\ast)\in\mathcal B
\]
mit
\[
J[\kappa^\ast,\phi^\ast]
=\min_{(\kappa,\phi)\in\mathcal B}J[\kappa,\phi].
\]
\end{definition}

\subsection{Interpretation}
Das Problem erweitert den klassischen Übergangsentwurf von reiner Geometrie
zum gekoppelten Geometrie--Überhöhungs-Entwurf und von Kurvenauswahl zur
funktionalen Optimierung. Der Übergangsbogen wird nicht mehr aus einem Katalog
gewählt, sondern entsteht als Teil eines optimalen gekoppelten Zustands.

\subsection{Sonderfälle}
\begin{proposition}
Für \(\phi(s)\equiv0\) reduziert sich das gekoppelte Problem auf reine
Krümmungsoptimierung.
\end{proposition}
\begin{proposition}
Gilt \(a_{\mathrm{eff}}(s)=0\) für alle \(s\), dann ist
\[
v^2\kappa(s)=g\sin\phi(s),
\]
und der Entwurf ist vollkommen ausgeglichen.
\end{proposition}
Der klassische Gleichgewichtsentwurf erscheint somit als Sonderfall.

\subsection{Bezug zur Optimalsteuerung}
Das gekoppelte Problem kann als Optimalsteuerungsproblem mit den Steuerungen
\[
\kappa(s),\quad\phi'(s),
\]
den Zustandsvariablen
\[
x(s),y(s),\theta(s),\phi(s)
\]
und dem Kostenfunktional \(J[\kappa,\phi]\) interpretiert werden. Damit entsteht
eine unmittelbare Verbindung zu modernen regelungstheoretischen und
variationsbasierten Verfahren.

\subsection{Verbindung zur Hauptthese}
\begin{summary}
Das gekoppelte Optimierungsproblem zeigt, dass Trassierungsentwurf keine rein
geometrische Aufgabe ist. Er optimiert ein gekoppeltes Zustandssystem, in dem
Krümmung und Überhöhung gemeinsam das dynamische Verhalten bestimmen.

Dies stützt die zentrale These, Übergangsbögen und Überhöhungsentwicklung als
Steuerungsfunktionen eines krümmungsbasierten Eisenbahnzustandsmodells zu
verstehen.
\end{summary}

\section{Zusammenfassung}
\begin{summary}
Alignment-Optimierung behandelt Geometrie als Variable statt als feste
Eingabe. Sie verbindet Parametrisierung, numerische Rekonstruktion und
beschränkte Optimierung, um geometrische und ingenieurtechnische Anforderungen
zu erfüllen.
\end{summary}
